\begin{statement}{A.24}{Proved}{Generic packet twists on the stable symmetric covers}
Let
\[
 C:\ v^2=w^7-w,\qquad
 D:\ z^2=(t^2+4)(t^3+3t),
\]
with the two cusps $P_0,P_\infty$ removed from $C$ and the single
cusp $Q_\infty$ removed from $D$. Use the maps
$C\to D\to E$ and the selected cusp paths of A.4. Fix a prime
$\ell$, any original packet section $s_b$ on $E\setminus O$,
and a rational prime $a\ne\ell$. Let
$\kappa_{a,\ell}:G_K\to\Z_\ell(1)$ be its Kummer cocycle, and
regard this as a cocycle in the $\ell$-component of
$I=\prod_p\Z_p(1)$. For $n\in\Z_\ell$ put
\[
 s_n(g)=c^{\,n\kappa_{a,\ell}(g)}s_b(g),
                                                        \tag{A.24.1}
\]
where all components of the displayed exponent other than the
$\ell$-component are zero.

There is one finite extension $L/K$, independent of $n$, such
that all these sections restrict to the selected cusp packets
on both covers over $G_L$. Put $\Delta_0=\Delta\cap G_L$.
Write $P_{C,\ell}$ and $P_{D,\ell}$ for the full pro-$\ell$
geometric groups of these two open curves, and $I_{C,\ell}$,
$I_{D,\ell}$ for their selected cusp inertia groups.

There is a set $B_{C,\ell}\subset\Z_\ell$ of cardinality at most
three such that, for every $n\notin B_{C,\ell}$ and every open
subgroup $\Delta_1\subset\Delta_0$,
\[
             P_{C,\ell}^{\,s_n(\Delta_1)}=I_{C,\ell}.
                                                        \tag{A.24.2}
\]
For $D$ the same assertion holds outside an at most countable
set $B_{D,\ell}$:
\[
             P_{D,\ell}^{\,s_n(\Delta_1)}=I_{D,\ell}.
                                                        \tag{A.24.3}
\]
At $\ell=2$ one may take $B_{D,2}=B_{C,2}$, so the bound is
again three. At every prime the complement of the indicated
exceptional set is dense in $\Z_\ell$ and has Haar measure one.
The sets are explicit in terms of the actual first nontrivial
Johnson matrices described in the proof. No value of those
arithmetic matrices, no membership decision for a specified
packet, and no assertion for all exceptional twists is made.
\end{statement}
\begin{proof}
\paragraph{A common arithmetic group and the factor $1/3$.}
The two maps to the base have ramification index three at the
selected cusps. For $C\to E_{0,-1}$, the coordinate map is
$(w^3,wv)$; its local parameter $x/y=w^2/v$ has order three
at $P_\infty$, and the map to $E$ is an unramified isogeny.
The involution $j$ giving $C\to D$ is fixed-point-free, so this
map is etale, also at the cusps, and the index for $D\to E$
is also three. With compatible generators the actual full
inertia inclusions are therefore
\[
                 c_C\longmapsto c^3,\qquad
                 c_D\longmapsto c^3.                    \tag{A.24.4}
\]
The image of $s_b$ lies in the arithmetic subgroup of the
chosen finite cover after restriction to a finite extension:
take the inverse image of that open subgroup under $s_b$.
The additional twist in (A.24.1) lies in the cover subgroup
automatically when $\ell\ne3$. When $\ell=3$, impose the one
further finite condition $\kappa_{a,3}=0\bmod3$.
Choose $L$ also to define the cusps, the tangent vectors, and
the isogenies $J_C\sim E^3$, $J_D\sim E^2$ of A.4.
These finitely many choices give a common $L$ for every $n$.
On $\Delta_0$ the induced twist in either cover is exactly
\[
                  c_C^{\,n\kappa_{a,\ell}/3}
             \quad\hbox{or}\quad
                  c_D^{\,n\kappa_{a,\ell}/3}.             \tag{A.24.5}
\]
At $\ell=3$ the numerator has been made divisible by three
in $\Z_3$; at all other primes three is a unit. Thus (A.24.5)
is an equality of integral pro-$\ell$ group actions, not
only a rational formula.

\paragraph{Tangential equivariance and trivial homology action.}
Let $h$ be the hyperelliptic involution of either curve.
It fixes the selected cusp, and its derivative there is $-1$.
There is a path in the punctured tangent line from a selected
tangent vector to its negative whose successive root
coordinates are roots of unity. All these coordinates are
fixed by $\Delta$, since $\mu_\infty\subset K^{\rm ab}$.
Use this path to make $h$ a based geometric automorphism.
Functoriality of the tangential fibre functor now gives an
actual commutation relation with the tangential section on
$\Delta_0$. This based automorphism fixes the chosen inertia
generator, so it also commutes with every twist of that
section by an inertia cocycle. It consequently commutes with
each packet action used here.

No assertion that this based lift has order two is needed.
Its square can be an inner automorphism by a cusp power.
Such an inner automorphism acts trivially on the
lower-central graded Lie algebra, and the induced action on
homology is the usual hyperelliptic one.

On $C$, that action is $+1$ on the one-dimensional boundary
space and $-1$ on proper homology. Over $\Q_\ell$ the two
eigenspaces split the homology extension equivariantly.
The isogeny $J_C\sim E^3$ makes the proper action of $\Delta_0$
trivial; the boundary action is the cyclotomic character,
also trivial on $\Delta_0$. Thus $\Delta_0$ acts trivially on
all of $H_1(C\setminus\{P_0,P_\infty\},\Z_\ell)$, since this
lattice is torsion-free. This argument does not divide an
integral lattice by two when $\ell=2$; the splitting is used
over $\Q_\ell$ only. For $D$, its one cusp contributes no
boundary homology, and $J_D\sim E^2$ gives the same
triviality directly. Every action in question is therefore
an IA automorphism of a free pro-$\ell$ group.

\paragraph{A nonzero Kummer value on $\Delta_0$.}
The restriction $\kappa_{a,\ell}|_\Delta$ is nonzero. To see
this directly, take a place $v$ of $K$ over $a$. Tame
$\ell$-inertia at $v$ has a generator $\tau$, and an
arithmetic Frobenius $F$ acts on it by the residue
cardinality $q_v$. The $\ell$-adic Kummer value of $a$ on
$\tau$ is its nonzero valuation $e_v$ times a chosen tame
generator. On the local commutator $[F,\tau]$, its value is
\[
                        (q_v-1)e_v\ne0
                       \quad\hbox{in }\Z_\ell(1).
                                                        \tag{A.24.6}
\]
Possible wild-inertia components do not affect this
calculation, because the residue characteristic is
$a\ne\ell$. The local commutator is in the global closed
commutator subgroup $\Delta$. On $\Delta$ the Kummer
cocycle is an additive character. A positive power of this
element belongs to the finite-index subgroup $\Delta_0$
and still has nonzero Kummer value. Fix such a
$\delta\in\Delta_0$, and put
\[
                         \beta=\kappa_{a,\ell}(\delta)/3
                              \in\Z_\ell\setminus\{0\}.
                                                        \tag{A.24.7}
\]
A compatible identification of $\Z_\ell(1)$ has been used
on $\Delta_0$ throughout.

\paragraph{The ordinary first Johnson map on $C$.}
Use the ordinary lower-central filtration $\gamma_d$ on
$P_{C,\ell}$. Its rational graded Lie algebra is
\[
                  \mathbb L(\Q_\ell t\oplus V),\qquad
                         \dim V=6,
\]
where $t$ is the class of $c_C$ and $V$ is the negative
hyperelliptic eigenspace. The class $t$ is primitive in
integral homology: a punctured genus-three surface with
two punctures has a free presentation in which either
selected cusp is a free generator.
The ordinary degree of $t$ is one. Its natural geometric
weight is two; these are different filtrations, and no
identification of them is used.

For the action of $s_b(\delta)$, the first Johnson map is
the degree-one derivation induced by
$g\mapsto s_b(\delta)g s_b(\delta)^{-1}g^{-1}$
modulo $\gamma_3$. It kills $t$, because the actual cusp
generator is fixed. Equivariance for $h$ forces it to send
$V$ into the negative eigenspace of
$\Lambda^2(\Q_\ell t\oplus V)$, namely $[t,V]$.
It therefore has the form
\[
             D_0(t)=0,\qquad D_0(v)=[t,A_\delta v],
                \qquad A_\delta\in\operatorname{End}_{\Q_\ell}(V).
\]
The integral twist (A.24.5) changes this to
\[
       D_n(t)=0,\qquad
       D_n(v)=[t,(A_\delta+n\beta\,1_V)v].               \tag{A.24.8}
\]
In geometric weights this part raises weight by two, which
explains why it appears in ordinary Johnson degree one.
There is also a useful constraint from the other cusp.
Let $\omega\in\Lambda^2V$ be its nondegenerate symplectic
class, written in the tensor algebra as
$\omega=\sum J_{ij}v_iv_j$ for a nonsingular skew matrix
$J$. We claim
\[
                         A_\delta J=J A_\delta^T.
                                                        \tag{A.24.13}
\]
For this calculation only, use the rational completed
free Lie algebra with $t=\log c_C$ as one generator.
The logarithm of a representative $c'$ of the other
cusp has the form
\[
               \log c'=-t-\omega+[t,u]
                         +\text{terms of degree at least three}.
\]
The term $[t,u]$ allows for the choice of the proper
homology complement and of the path to that cusp.
The packet action preserves the conjugacy class of
$c'$ and its exponent, since $\chi_\ell(\delta)=1$.
If $f_1$ is the linear term of a conjugator, comparison
in degree two gives $[f_1,t]=0$, hence $f_1=a t$.
Comparison in degree three, using that the action
fixes the actual $t$, now gives
$D_0(\omega)\in[t,\mathbb L_2(\Q_\ell t\oplus V)]$.
The extra term $D_0([t,u])$ is itself in that subspace.
Project to words containing one $t$ and two letters of
$V$. The words $v_i t v_j$ therefore have coefficient
zero. Their coefficient matrix in $D_0(\omega)$ is
$-A_\delta J+J A_\delta^T$, which proves (A.24.13).
This is the peripheral constraint on the actual action,
not a condition imposed on a freely chosen derivation.

It follows that $(A_\delta+n\beta\,1_V)J$ is skew.
The polynomial
\[
 p_\delta(n)=
 \frac{\operatorname{Pf}((A_\delta+n\beta\,1_V)J)}
      {\operatorname{Pf}(J)}
 \quad\hbox{satisfies}\quad
 \det(A_\delta+n\beta\,1_V)=p_\delta(n)^2.
                                                        \tag{A.24.14}
\]
Here $\operatorname{Pf}$ is the Pfaffian of a six-by-six
skew matrix. It has degree three in its entries, so
$p_\delta$ has degree three in $n$ and leading
coefficient $\beta^3\ne0$.
Let $B_{C,\ell}$ be its roots in $\Z_\ell$.

We prove the Lie-kernel assertion needed here in all
degrees. For a vector space $V$ over a characteristic-zero
field and an invertible $A$, let
$D(t)=0$, $D(v)=[t,Av]$. Then
\[
                    \ker D=\Q_\ell t.                  \tag{A.24.9}
\]
Indeed let $V_j=\operatorname{ad}_t^jV$ for $j\ge0$.
The ideal generated by $V$ is freely generated by
$W=\bigoplus_{j\ge0}V_j$. One can verify this elimination
identity without an additional relation theorem:
on $\mathbb L(W)$ the shift $V_j\to V_{j+1}$ is a
derivation; adjoining $t$ acting by that shift gives the
universal Lie algebra on $t$ and $V_0$.
The derivation under consideration satisfies
$D(v_j)=(Av)_{j+1}$.
Embed $\mathbb L(W)$ in the free associative algebra.
For a nonzero polynomial choose its largest occurring jet
index $j$. In its derivative, the terms containing a new
letter from $V_{j+1}$ come only from replacing a letter in
$V_j$. If those terms vanished, substitute
$V_{j+1}\to A^{-1}V_j$. The result is the Euler operator
counting the occurrences of $V_j$ in each word.
It cannot annihilate the part involving $V_j$ in
characteristic zero. Thus the derivative is nonzero on the
whole ideal. This proves (A.24.9).

If $u\ne1$ is fixed by $s_n(\delta)$, its first nonzero
lower-central class is killed by $D_n$. When
$n\notin B_{C,\ell}$, (A.24.9) forces that class to be
$a t$ in degree one. Its coefficient is in $\Z_\ell$,
because $t$ is primitive in the integral lattice.
Subtract the actual fixed cusp power $c_C^a$.
Any nonidentity remainder would have first degree at least
two, contradicting (A.24.9). The filtration is separated,
so $u=c_C^a$. This is an argument about actual pro-$\ell$
elements, with no lifting of a rational Malcev element.
It proves (A.24.2) for $\Delta_0$.
For any open $\Delta_1$, take a positive power
$\delta^N\in\Delta_1$; its first Johnson map is $ND_n$,
which has the same kernel. This proves the assertion for
every stated open subgroup.

\paragraph{The ordinary degree-two pencil on $D$.}
Now $P_{D,\ell}$ is free pro-$\ell$ of rank four. Write
$W=H_1(D\setminus Q_\infty,\Q_\ell)$ and choose a symplectic
integral basis so that the initial cusp class is
\[
                 \omega=[x_1,y_1]+[x_2,y_2]
                      \in\operatorname{gr}_2P_{D,\ell}.
\]
This class is primitive over $\Z_\ell$, and the cusp has
no degree-one term. The hyperelliptic involution acts as
$-1$ on $W$ and as $+1$ on $\mathbb L_2(W)$.
Its commutation with the packet action therefore forces
the degree-one Johnson map to vanish. The degree-two map
for $s_n(\delta)$ is consequently a well-defined
homogeneous derivation of the form
\[
          T_n=T_0+n\beta\,\operatorname{ad}(\omega),
                         \qquad T_n(\omega)=0.           \tag{A.24.10}
\]
This formula uses the actual first degree-two term of the
action and asserts no formula for $T_0$.

For clarity, the required centralizer fact has a short
direct proof:
\[
                C_{\mathbb L(W)}(\omega)=\Q_\ell\omega.
                                                        \tag{A.24.11}
\]
In the tensor algebra order the letters so that the
leading word of $\omega$ is $x_1y_1$.
This word is not a proper power. If a homogeneous
associative polynomial $f$ commutes with $\omega$, their
leading words commute, so the leading word of $f$ is
$(x_1y_1)^m$. Subtracting the matching multiple of
$\omega^m$ would leave a commuting polynomial with a
strictly smaller leading word, which is impossible by
the same word argument. Thus $f$ is a scalar multiple of
$\omega^m$. A power $\omega^m$ with $m>1$ is not primitive
for the tensor Hopf coproduct in characteristic zero,
whereas every Lie element is primitive. This proves
(A.24.11).

In each degree $d\ne2$,
$\operatorname{ad}(\omega):\mathbb L_d(W)\to
\mathbb L_{d+2}(W)$ is therefore injective. In degree two
use the source $\mathbb L_2(W)/\Q_\ell\omega$.
For each of these finite-dimensional source spaces
choose a full-rank minor which is nonzero for
$\operatorname{ad}(\omega)$.
The corresponding minor for (A.24.10) is a polynomial
in $n$ with nonzero leading coefficient. Let $B_{D,\ell}$
be the union, over all $d\ge1$, of its roots in $\Z_\ell$.
This set is at most countable. Outside it,
\[
                         \ker T_n=\Q_\ell\omega.
                                                        \tag{A.24.12}
\]
The first nonzero class of a fixed pro-$\ell$ element is
then $a\omega$ in degree two. Primitivity makes
$a\in\Z_\ell$; subtract the actual cusp power $c_D^a$.
A nonidentity remainder has first degree at least three
and contradicts (A.24.12). This proves (A.24.3).
Replacing $\delta$ by a positive power proves it for
all open $\Delta_1$ just as above.

\paragraph{The sharper bound at two and the scope.}
For $\ell=2$ the etale double cover $C\to D$ induces an
embedding of $P_{C,2}$ as an open subgroup of $P_{D,2}$,
by the pro-prime topology argument in A.22.
If $u\in P_{D,2}$ is fixed, then $u^2\in P_{C,2}$ is
fixed. Outside $B_{C,2}$ it belongs to the selected
$I_{C,2}$, whose image is $I_{D,2}$ since this double
cover is unramified at the cusp.
The latter inertia is root-closed, as follows here
directly from (A.24.11). If $v$ commutes with
$c_D^b\ne1$, its first lower-central class commutes with
$b\omega$. It is therefore $a\omega$ in degree two,
with $a\in\Z_2$ by primitivity. Subtracting $c_D^a$
leaves an element still commuting with $c_D^b$ and
having first degree at least three, which (A.24.11)
excludes. Hence $C_{P_{D,2}}(c_D^b)=I_{D,2}$.
If $u^2=c_D^b\ne1$, apply this to $u$ itself.
If $u^2=1$, any first nonzero lower-central class of
$u$ would be killed by two in a torsion-free
$\Z_2$-module, so $u=1$ as well.
This proves (A.24.3) outside $B_{C,2}$.

Finite and countable subsets of $\Z_\ell$ have Haar
measure zero, and their complements are dense because
every nonempty open ball is uncountable.
The independent certificate
\texttt{code/a\_generic\_packets.py} checks the tangent
signs, the ramification orders, and finite-degree Lie
ranks for deliberately resonant invertible matrices.
It does not compute $A_\delta$ or $T_0$ and is not used
to substitute for the unbounded Lie proofs above.
\end{proof}
